\rtmtight
\begin{tblr}{width=\textwidth,colspec={|Q[c,m,2.1cm]|X[l,m]|},hlines,vlines,hspan=minimal,rowsep=1.5pt,colsep=3pt,row{1}={bg=headgray,font=\bfseries}}
\SetCell{c,m}\hfil\logo\hfil & \textbf{POLITEKNIK NEGERI MALANG}\par\textbf{JURUSAN TEKNOLOGI INFORMASI}\par\textbf{PROGRAM STUDI: D4 TEKNIK INFORMATIKA} \\
\SetCell[c=2]{l} \textbf{RENCANA TUGAS MAHASISWA (RTM) DAN RUBRIK PENILAIAN} & \\
Nama Mata Kuliah & Basis Data Lanjut \\
Kode Mata Kuliah & RTI253005 \quad \textbf{sks / Jam perminggu:} 2 SKS/4 jam \quad \textbf{Semester:} 3 \\
Dosen Pengampu & \begin{enumerate}\item Candra Bella Vista, S.Kom., M.T.\item M. Hasyim Ratsanjani, S.Kom., M.Kom.\item Ridwan Rismanto, S.ST., M.Kom., Ph.D.\end{enumerate} \\
\end{tblr}
\assessmentblock{Tugas Jobsheet Praktikum PostgreSQL}{Tugas terstruktur berbasis praktikum (PBL)}{SCPMK0203-01601, SCPMK0203-01602, SCPMK0701-01603, SCPMK1001-01604}{Mahasiswa menyelesaikan serangkaian jobsheet praktikum PostgreSQL mencakup query dasar, query lanjutan, indeks, objek basis data, transaksi, FTS, JSONB, backup-restore, migrasi, dan integrasi dengan aplikasi.}{Menganalisis studi kasus, mengeksekusi perintah SQL/PostgreSQL sesuai jobsheet, mendokumentasikan hasil eksekusi beserta penjelasan, dan mengumpulkan melalui LMS.}{Skrip SQL, screenshot hasil eksekusi, laporan singkat per jobsheet dalam format PDF atau dokumen digital.}{Ketepatan sintaks dan keberhasilan eksekusi query & 10 & Semua perintah SQL/PostgreSQL benar secara sintaks, berhasil dieksekusi tanpa error, dan menghasilkan output sesuai kebutuhan studi kasus. & Sebagian besar perintah benar dan berhasil dieksekusi, tetapi ada beberapa error minor atau output yang kurang tepat. & Banyak perintah salah sintaks, gagal dieksekusi, atau output tidak sesuai kebutuhan studi kasus. \\ \\
Kualitas analisis dan pemecahan masalah & 8 & Mahasiswa menganalisis studi kasus dengan tepat, memilih pendekatan yang optimal, dan menjelaskan alasan pemilihan teknik atau perintah yang digunakan. & Analisis studi kasus cukup tepat, tetapi penjelasan alasan pemilihan teknik atau optimasi masih terbatas. & Analisis studi kasus tidak tepat, tidak ada penjelasan, atau pendekatan yang dipilih tidak sesuai kebutuhan. \\ \\
Kelengkapan dokumentasi dan laporan & 7 & Semua langkah terdokumentasi dengan lengkap, screenshot hasil eksekusi tersedia, dan laporan menjelaskan proses secara runtut dan mudah dipahami. & Dokumentasi sebagian besar lengkap, tetapi beberapa langkah atau screenshot masih kurang atau tidak jelas. & Dokumentasi tidak lengkap, screenshot tidak tersedia, atau laporan tidak menjelaskan proses pengerjaan. \\}

\assessmentblock{Kuis 1 SQL Dasar dan Lanjutan}{Kuis (ujian tertulis/praktik PostgreSQL, closed book)}{SCPMK0203-01601}{Mahasiswa mengerjakan Kuis 1 materi minggu 1--3 meliputi pengenalan PostgreSQL, DDL, DML, sub-query, JOIN, grouping, dan CTE secara mandiri.}{Mengerjakan soal berbasis kasus SQL menggunakan PostgreSQL secara individu dengan sifat closed book sesuai jadwal asesmen.}{Lembar jawaban atau file skrip SQL hasil kuis.}{Penguasaan DDL dan DML PostgreSQL & 4 & Perintah CREATE, ALTER, DROP, INSERT, UPDATE, DELETE benar sintaks dan sesuai kebutuhan studi kasus dengan constraint yang tepat. & Sebagian besar perintah DDL/DML benar, tetapi ada kesalahan pada constraint atau tipe data. & Banyak perintah DDL/DML salah sintaks atau tidak sesuai kebutuhan studi kasus. \\ \\
Ketepatan query lanjutan (sub-query dan JOIN) & 4 & Sub-query skalar/korelasi dan semua jenis JOIN digunakan dengan benar dan menghasilkan data yang tepat sesuai kebutuhan kasus. & Sebagian besar query sub-query dan JOIN benar, tetapi ada kesalahan pada kondisi join atau hasil yang kurang tepat. & Query sub-query atau JOIN salah secara logika atau tidak dapat dieksekusi. \\ \\
Penerapan CTE dan grouping & 2 & CTE dan GROUP BY dengan HAVING digunakan secara tepat untuk menyederhanakan query kompleks dan menghasilkan agregasi yang benar. & CTE atau grouping digunakan tetapi ada kesalahan pada struktur atau kondisi HAVING. & CTE atau grouping tidak digunakan atau menghasilkan hasil yang salah. \\}

\assessmentblock{UTS Query dan Objek Basis Data PostgreSQL}{UTS (ujian praktik PostgreSQL, closed book)}{SCPMK0203-01601, SCPMK0203-01602, SCPMK0701-01603}{Mahasiswa mengerjakan soal UTS materi minggu 1--8 meliputi query dasar, query lanjutan, indeks, function, view, stored procedure, transaksi, FTS, dan JSONB secara mandiri.}{Mengerjakan soal berbasis studi kasus PostgreSQL secara individu dengan sifat closed book sesuai jadwal asesmen.}{File skrip SQL dan/atau lembar jawaban UTS.}{Ketepatan query dasar dan lanjutan & 6 & Query DDL, DML, sub-query, JOIN, dan CTE benar dan menghasilkan data sesuai studi kasus UTS. & Sebagian besar query benar, tetapi ada kesalahan logika atau output yang kurang tepat. & Banyak query salah atau tidak dapat dieksekusi. \\ \\
Ketepatan implementasi indeks dan objek basis data & 7 & Indeks dipilih dan dibuat dengan tepat; function, view, dan stored procedure berfungsi benar sesuai spesifikasi studi kasus. & Implementasi indeks dan objek basis data sebagian besar benar, tetapi ada ketidaktepatan pada logika atau tipe indeks. & Indeks atau objek basis data tidak dapat dibuat atau tidak berfungsi sesuai kebutuhan. \\ \\
Penerapan transaksi, FTS, dan JSONB & 7 & Transaksi ACID, full-text search, dan query JSONB diimplementasikan benar sesuai kebutuhan studi kasus. & Implementasi transaksi, FTS, atau JSONB sebagian besar benar, tetapi ada kasus yang tidak tertangani. & Transaksi tidak konsisten, FTS tidak berfungsi, atau query JSONB salah secara logika. \\}

\assessmentblock{PBL Proyek Basis Data Terintegrasi}{Project Based Learning}{SCPMK1001-01604}{Mahasiswa secara berkelompok merancang, mengimplementasikan, dan mempresentasikan proyek basis data PostgreSQL terintegrasi berdasarkan domain nyata yang dipilih, mencakup perancangan skema, optimasi, stored procedure, transaksi, dan integrasi dengan aplikasi.}{Memilih domain proyek, menganalisis kebutuhan data, merancang skema, mengimplementasikan seluruh fitur PostgreSQL yang relevan, mengintegrasikan dengan aplikasi, menguji, dan mempresentasikan hasilnya.}{Laporan proyek, skrip SQL lengkap, kode aplikasi, dokumentasi pengujian, dan bahan presentasi.}{Kualitas rancangan dan implementasi basis data & 9 & Skema dirancang dengan baik (normalisasi, constraint, indeks), implementasi mencakup seluruh fitur yang direncanakan, dan tidak ada error signifikan. & Skema dan implementasi sebagian besar baik, tetapi ada kekurangan pada normalisasi, indeks, atau beberapa fitur belum terimplementasi. & Skema tidak terancang dengan baik atau banyak fitur yang direncanakan tidak terimplementasi. \\ \\
Integrasi dengan aplikasi dan fungsionalitas & 7 & Basis data terintegrasi penuh dengan aplikasi, semua operasi CRUD berfungsi, transaksi dikelola dengan benar, dan aplikasi berjalan tanpa error. & Integrasi berjalan untuk fungsi utama, tetapi ada fitur yang tidak berfungsi atau penanganan error yang belum lengkap. & Integrasi dengan aplikasi tidak berfungsi atau hanya sebagian kecil fitur yang berjalan. \\ \\
Kualitas presentasi dan argumentasi teknis & 4 & Presentasi terstruktur, alasan teknis desain disampaikan dengan jelas, dan tim mampu menjawab pertanyaan dengan baik. & Presentasi cukup jelas, tetapi argumentasi teknis atau jawaban pertanyaan masih kurang mendalam. & Presentasi tidak terstruktur atau tim tidak dapat menjelaskan keputusan teknis yang diambil. \\}

\assessmentblock{UAS Proyek Akhir Basis Data PostgreSQL}{UAS (defense dan demonstrasi proyek)}{SCPMK1001-01604}{Mahasiswa mempresentasikan dan mendemonstrasikan proyek basis data PostgreSQL terintegrasi yang mencakup seluruh kompetensi mata kuliah, disertai defense atas keputusan teknis yang diambil.}{Mempersiapkan demonstrasi live proyek, menyusun bahan presentasi yang komprehensif, dan menjawab pertanyaan penguji terkait aspek teknis dan arsitektur basis data.}{Demonstrasi live aplikasi, skrip SQL final, laporan proyek lengkap, dan bahan presentasi.}{Kelengkapan fitur dan kualitas implementasi proyek & 10 & Semua fitur yang direncanakan terimplementasi dengan baik, performa basis data dioptimalkan, dan aplikasi berjalan stabil tanpa error. & Sebagian besar fitur terimplementasi, tetapi ada fitur yang belum optimal atau ada error yang masih muncul dalam kondisi tertentu. & Banyak fitur yang tidak terimplementasi atau aplikasi tidak dapat berjalan dengan baik saat demonstrasi. \\ \\
Kualitas arsitektur dan keputusan teknis & 9 & Arsitektur basis data dirancang secara sistematis, keputusan teknis (indeks, transaksi, objek basis data) dapat dipertahankan dengan argumen yang kuat. & Arsitektur cukup baik, tetapi beberapa keputusan teknis tidak optimal atau argumen yang diberikan kurang kuat. & Arsitektur tidak terancang dengan baik atau mahasiswa tidak dapat menjelaskan keputusan teknis yang diambil. \\ \\
Kemampuan defense dan komunikasi teknis & 6 & Mahasiswa menjawab pertanyaan penguji dengan tepat, menjelaskan konsep teknis dengan bahasa yang jelas, dan mampu menghubungkan implementasi dengan teori. & Mahasiswa menjawab sebagian besar pertanyaan dengan baik, tetapi ada beberapa pertanyaan teknis yang tidak dapat dijawab secara mendalam. & Mahasiswa tidak dapat menjawab pertanyaan teknis atau penjelasan tidak menunjukkan pemahaman yang cukup. \\}

\begin{tblr}{width=\textwidth,colspec={|Q[l,m,3.2cm]|X[l,m]|},hlines,vlines,hspan=minimal,row{1-3}={bg=headgray},rowsep=1.5pt,colsep=3pt}
Jadwal Pelaksanaan: & Minggu 1--17 sesuai RPS. Setiap evaluasi memiliki RTM dan rubrik multi-kriteria. \\
Lain-Lain Yang Diperlukan: & LMS, PostgreSQL, pgAdmin/DBeaver, repositori kode, dan perangkat presentasi. \\
Pustaka: & Mengacu pustaka utama dan pendukung pada bagian RPS. \\
\end{tblr}
